% !TeX root = ../../Book1.tex
\OptionalChapter{测度的弱收敛与紧性}{b1:ch:12}

% 资料组：Billingsley1999Convergence 为已核验出版信息/目录的后续伴读，
% Yan2004Cedulun 第二版第6.2--6.5节为中文目录与名称对照，不承担未读证明；
% Fremlin2016Measure 卷四437J--P、437U--V用于拓扑约定与一般紧性版本。
% Fremlin 的 vague 使用 Cb 测试，且原书 Radon 约定涉及完备化；
% 本书淡收敛固定 Cc 测试，弱收敛固定 Cb 测试，测度先定义于 Borel 集。
% 证明义务：测试空间比较、集合边界与弱极限、共同紧集控制、概率弱拓扑度量化、
% 紧立方体中的极限拉回、Radon与Prokhorov条件以及拓扑/序列紧性的边界。
\begin{chapterabstract}
上一章用连续函数上的积分恢复测度，本章进一步让测度本身随序列变化。
只观察每个固定紧区域，可能看不到逐渐远去的质量；把测试范围扩大到所有有界连续函数，
则连总质量也必须参与极限。我们据此区别淡收敛与弱收敛，
用 Portmanteau 定理把连续测试改写成开集、闭集及边界零测集上的判据，
再用紧族条件控制整列测度的尾部。
在波兰空间上，Prokhorov 定理说明这种共同的紧集控制恰好对应相对弱紧性，
从而提供抽取收敛子列的方法。极限的存在与辨认仍是两项不同的工作：
前者依赖紧性，后者依赖测试函数、边缘分布或其他具体结构。
\end{chapterabstract}

\begin{learninggoals}
\begin{itemize}
\item\label{b1:item:12-goal-tests}区别 \(C_c\)、\(C_0\) 与 \(C_b\) 测试，以及测度弱收敛和 Banach 空间弱拓扑的不同含义。
\item\label{b1:item:12-goal-portmanteau}由连续测试推导开闭集不等式，正确使用边界零测集与几乎处处连续测试的判据。
\item\label{b1:item:12-goal-tightness}构造对整族有效的紧集，说明淡收敛在何种条件下保留总质量并提升为弱收敛。
\item\label{b1:item:12-goal-prokhorov}掌握波兰空间上的 Prokhorov 定理，分清 Radon 正则性、相对紧性与收敛子列各自所需的条件。
\item\label{b1:item:12-goal-applications}用矩估计、有限原子逼近和具体函数空间中的紧集处理分布极限，并检验无界测试的尾部条件。
\end{itemize}
\end{learninggoals}

本章需要第\ref{b1:ch:05}章的积分极限定理、第\ref{b1:ch:09}章的弱-*紧性与度量化思想，
以及第\ref{b1:ch:11}章的连续截断和 Riesz--Markov 表示。
乘积测度习题还会调用第\ref{b1:ch:07}章的 Fubini 定理。
概率测度在这里首先只是总质量为 \(1\) 的正测度；谈论分布时使用可测映射的推前，
不预设随机过程课程。拓扑方面需要熟悉紧集、可分性及完备度量，
本章所需的有限测度正则性和概率测度空间的度量化均在正文中建立。

首次阅读可先以实轴上的概率分布为模型。
第\ref{b1:sec:1201}节决定哪些测试参与收敛；第\ref{b1:sec:1202}节回答这些测试怎样控制集合的质量。
随后在第\ref{b1:sec:1203}节观察统一截断的作用，特别留意“每个测度都紧”与“整族是紧族”的量词差别。
第\ref{b1:sec:1204}节不再要求局部紧性，而是在波兰空间上通过紧立方体与测度表示取得聚点，
再把它拉回原空间。最后一个小节的一般拓扑版本可以留待第二遍阅读；
其中保留的是 Radon 概率测度族的拓扑紧性，不能自动改写为序列结论。

希望继续学习概率分布的极限理论，可将 Patrick Billingsley 的《Convergence of Probability Measures》%
\cite[第二版]{Billingsley1999Convergence}作为后续专书。
中文伴读可选严加安的《测度论讲义》第二版\cite[第6.2--6.5节]{Yan2004Cedulun}，
按其目录对照弱收敛、测度族的紧性与淡收敛在中文教材中的组织次序。
这里的两项推荐用于阅读导航，不替代本章给出的证明。

较一般的拓扑测度空间可参阅 D.~H. Fremlin 的《Measure Theory》%
\cite[卷四，437J--P、437U--V]{Fremlin2016Measure}。
该书区分不同的测度拓扑并讨论 Prokhorov 空间，适合核对一般版本的条件。
阅读时须先对齐约定：该书的 vague 使用 \(C_b\) 测试，本书的淡收敛则使用 \(C_c\)；
其 Radon 测度的完备性约定也不同于本书先在 Borel 集上讨论测度的方式。
后文引用相应结论时，均以具体测试类和测度范围为准。

% !TeX root = ../../Book1.tex
\section{淡收敛与弱收敛}
\label{b1:sec:1201}

把一个测度看成积分泛函以后，测度列的收敛便可以通过积分值来定义。
不过，测试函数的范围不能省略。同一列概率测度，用紧支集函数测试可能趋于零，
用常函数测试却始终得到一。这里的差别不是记号上的习惯，而是是否保留整体质量。

\ProseParagraphBreak

本节先在局部紧 Hausdorff 空间上接续前一章的表示定理，再转到度量空间上的概率测度。
本章未特别说明时，测度均为正 Borel 测度；需要符号或复测度时会另行指出。

\subsection{测试函数}
\label{b1:subsec:120101}

对拓扑空间 \(X\)，记
\[
 C_b(X)\coloneqq\{\,f:X\to\mathbb R:f\text{ 连续且有界}\,\},
 \quad \|f\|_\infty\coloneqq\sup_{x\in X}|f(x)|.
\]
\symidx{\(C_b(X)\)：有界连续函数空间}
它在一致范数下是 Banach 空间。一致 Cauchy 列的逐点极限仍有界，且收敛一致，
因而极限连续。以下用实值函数测试已经足够；若要使用复值函数，分别测试实部和虚部即可。

\ProseParagraphBreak

当 \(X\) 局部紧 Hausdorff 时，有
\[
 C_c(X)\subset C_0(X)\subset C_b(X).
\]
三个空间分别允许我们观察紧区域、无穷远处逐渐减弱的区域，以及整个空间。
若 \(X\) 紧，三者相同；若 \(X\) 非紧，常函数 \(1\) 属于 \(C_b(X)\)，
却不属于 \(C_0(X)\)。对局部有限的测度，\(C_c(X)\) 上的积分总是有限的；
要让全部有界测试函数都可积，则必须要求测度有限。

\ProseParagraphBreak

有限符号或复 Radon 测度与 \(C_0(X)^*\) 的等距对应已在%
\cref{b1:thm:riesz-markov-c0}中建立。它给出了一种现成的弱*\WeakStarJoin{}收敛：
对每个 \(f\in C_0(X)\)，积分 \(\int f\,\mathrm{d}\mu_n\) 收敛。
但本章还要使用 \(C_c(X)\) 和 \(C_b(X)\)。每当比较两种收敛时，
应先写出测试空间，再判断哪个包含关系或逼近结论能够发挥作用。

\subsection{淡收敛}
\label{b1:subsec:120102}

\begin{definition}\label{b1:def:vague-convergence}
设 \(X\) 为局部紧 Hausdorff 空间，\(\mu_n,\mu\) 为其上的 Radon 测度。
若对每个 \(f\in C_c(X)\) 都有
\[
 \lim_{n\to\infty}\int_X f\,\mathrm{d}\mu_n=\int_X f\,\mathrm{d}\mu,
\]
则称 \(\mu_n\) \term[danshoulian]{淡收敛}[vague convergence]于 \(\mu\)，
记为 \(\mu_n\xrightarrow{v}\mu\)。
\end{definition}
\symidx{\(\mu_n\xrightarrow{v}\mu\)：由连续紧支集函数测试的淡收敛}

中文名称采用严加安的《测度论讲义》\cite[第6.5节]{Yan2004Cedulun}。
不同文献对 vague 的测试空间并不一致；例如 Fremlin 的《Measure Theory》%
\cite[437J(c)]{Fremlin2016Measure}使用 \(C_b(X)\) 定义该名称。
本书始终按上面的 \(C_c(X)\) 约定，不凭相同的外文名称移用结论。

\ProseParagraphBreak

淡极限若存在便唯一：两个候选极限在全部 \(C_c(X)\) 上有相同积分，
由正测度表示的唯一性即相等。它还自动控制每一块固定紧集上的质量。

\begin{proposition}\label{b1:prop:vague-local-bound}
若 \(\mu_n\xrightarrow{v}\mu\)，则对每个紧集 \(K\subset X\)，
有 \(\sup_n\mu_n(K)<\infty\)。
\end{proposition}
\begin{proof}
由\cref{b1:lem:lch-cutoff}取 \(h\in C_c(X)\)，使 \(0\leq h\leq1\)、
\(h|_K=1\)。于是
\[
 \mu_n(K)\leq\int_X h\,\mathrm{d}\mu_n\longrightarrow\int_X h\,\mathrm{d}\mu<\infty.
\]
收敛数列有界，结论随之成立。
\end{proof}

此处的界依赖 \(K\)，不能直接对所有紧集取上确界。例如在 \(\mathbb R\) 上令
\(\mu_n=n\delta_n\)。每个紧支集测试函数最终在点 \(n\) 处为零，
所以 \(\mu_n\xrightarrow{v}0\)，但 \(\mu_n(\mathbb R)=n\)。
甚至 \(C_0(\mathbb R)\) 中的函数 \(f(x)=(1+|x|)^{-1/2}\) 也满足
\(\int f\,\mathrm{d}\mu_n=n/\sqrt{1+n}\to\infty\)。
因此没有整体界时，不能把紧支集测试悄悄改为无穷远消失的测试。

\begin{proposition}\label{b1:prop:vague-c0-bounded}
设 \(\mu_n,\mu\) 为局部紧 Hausdorff 空间上的有限正 Radon 测度，并且
\(M=\sup_n\mu_n(X)<\infty\)。则 \(\mu_n\xrightarrow{v}\mu\) 当且仅当
它们在 \(C_0(X)^*\) 中弱*\WeakStarJoin{}收敛于 \(\mu\)。
对有限符号或复 Radon 测度，若改用 \(\sup_n|\mu_n|(X)<\infty\)，结论仍成立。
\end{proposition}
\begin{proof}
弱*\WeakStarJoin{}收敛蕴含淡收敛，因为 \(C_c(X)\subset C_0(X)\)。反过来，给定
\(f\in C_0(X)\) 和 \(\varepsilon>0\)，由\cref{b1:prop:cc-dense-c0}%
选 \(g\in C_c(X)\)，使 \(\|f-g\|_\infty<\varepsilon\)。于是
\[
 \left|\int_X f\,\mathrm{d}\mu_n-\int_X f\,\mathrm{d}\mu\right|
 \leq\left|\int_X g\,\mathrm{d}\mu_n-\int_X g\,\mathrm{d}\mu\right|
       +\varepsilon\bigl(M+\mu(X)\bigr).
\]
先令 \(n\to\infty\)，再令 \(\varepsilon\downarrow0\)，便得到所需收敛。
一般标量测度的证明相同，把两个质量替换为全变差质量即可。
\end{proof}

上式中的统一界有明确用途：同一个一致逼近误差，必须对所有 \(n\) 同时有效。
前一例正是在这一步失去控制。另一方面，若已知在 \(C_0(X)^*\) 中弱*\WeakStarJoin{}收敛，
一致有界原理又保证范数一致有界，因此对于收敛的序列，弱*\WeakStarJoin{}条件本身包含了全变差界。

\subsection{弱收敛}
\label{b1:subsec:120103}

\begin{definition}\label{b1:def:measure-weak-convergence}
设 \(X\) 为拓扑空间，\(\mu_n,\mu\) 为其上的有限正 Borel 测度。
若对每个 \(f\in C_b(X)\) 都有
\[
 \lim_{n\to\infty}\int_X f\,\mathrm{d}\mu_n=\int_X f\,\mathrm{d}\mu,
\]
则称 \(\mu_n\) \term[cedu-ruoshoulian]{弱收敛}[weak convergence of measures]于 \(\mu\)，
记为 \(\mu_n\Rightarrow\mu\)。测度空间上使所有映射
\(\nu\mapsto\int f\,\mathrm{d}\nu\)、\(f\in C_b(X)\) 连续的最弱拓扑，称为这里的弱拓扑。
\end{definition}
\symidx{\(\mu_n\Rightarrow\mu\)：由有界连续函数测试的测度弱收敛}

这个拓扑在 \(\mu\) 处的基本邻域可写成
\[
 \left\{\nu:
 \left|\int_X f_j\,\mathrm{d}\nu-\int_X f_j\,\mathrm{d}\mu\right|<\varepsilon,
 \quad 1\leq j\leq m\right\},
\]
其中只选有限多个测试函数。对于一般拓扑空间，连续函数未必足以区别测度；
后文主要使用度量空间，或局部紧空间上的 Radon 测度，这两个场合没有这样的歧义。

\begin{proposition}\label{b1:prop:weak-mass-uniqueness}
若 \(\mu_n\Rightarrow\mu\)，则 \(\mu_n(X)\to\mu(X)\)。
在度量空间上，有限正 Borel 测度由其在 \(C_b(X)\) 上的积分唯一确定；
在局部紧 Hausdorff 空间上，有限正 Radon 测度也具有此性质。
\end{proposition}
\begin{proof}
第一点取 \(f=1\)。在度量空间上，对非空闭集 \(F\) 定义
\[
 h_k(x)=\max\{1-k\cdot \mathrm{d}(x,F),0\}.
\]
这些函数有界连续，且 \(h_k\downarrow\mathbf1_F\)。若两个有限测度的连续测试积分相同，
由支配收敛定理，它们在 \(F\) 上的值也相同。空集的值当然相同。
所有闭集构成包含 \(X\) 的 \(\pi\)-系并生成 Borel \(\sigma\)-代数，
故有限测度的唯一性定理给出两测度相同。局部紧情形则直接使用
\(C_c(X)\subset C_b(X)\) 及前一章的 Radon 表示唯一性。
\end{proof}

这里的测度弱收敛，不是把测度 Banach 空间 \(M(X)=C_0(X)^*\) 的整个对偶再拿来测试。
后者对应 \(\sigma\bigl(M(X),M(X)^*\bigr)\)，是\chapref{b1:ch:09}意义下 Banach 空间的弱拓扑。
本章固定 \(C_b(X)\) 测试，不能仅因名称中都有“弱”字便将它们等同。
即使 \(X\) 紧，测度弱拓扑对应的也是 \(C(X)^*\) 的弱*\WeakStarJoin{}拓扑，而非一般的弱拓扑。

\subsection{概率测度情形}
\label{b1:subsec:120104}

记 \(\mathcal P(X)\) 为 \(X\) 上所有 Borel 概率测度构成的集合，赋予刚定义的弱拓扑。
\symidx{\(\mathcal P(X)\)：Borel 概率测度空间}
若有限正测度列的每一项都有质量一，那么任何弱极限仍有质量一。
但淡极限不一定如此；缺少的质量可能已经离开每一个固定的紧区域。
\secref{b1:sec:1203}将用紧族把这种现象精确表述出来。

\ProseParagraphBreak

先看与质量逃逸不同的集中现象。在 \(\mathbb R\) 上令
\(\mathrm{d}\mu_n(x)=n\cdot\mathbf1_{[0,1/n]}(x)\,\mathrm{d}x\)。对任意 \(f\in C_b(\mathbb R)\)，有
\[
 \left|\int_{\mathbb R}f\,\mathrm{d}\mu_n-f(0)\right|
 \leq\sup_{0\leq x\leq1/n}|f(x)-f(0)|\longrightarrow0.
\]
所以 \(\mu_n\Rightarrow\delta_0\)。每个 \(\mu_n\) 都由密度给出，极限却是点质量。
此外，\(\mu_n(\{0\})=0\) 而 \(\delta_0(\{0\})=1\)，说明弱收敛不要求逐个
Borel 集的测度值收敛。下一节会指出哪些集合仍能保留这一结论。

\begin{proposition}[连续映射原理]\label{b1:prop:weak-continuous-mapping}
设 \(T:X\to Y\) 连续，\(\mu_n,\mu\) 为 \(X\) 上的有限正 Borel 测度。
若 \(\mu_n\Rightarrow\mu\)，则推前测度满足 \(T_\#\mu_n\Rightarrow T_\#\mu\)。
事实上，\(\mu\mapsto T_\#\mu\) 关于两边的弱拓扑连续。
\end{proposition}
\begin{proof}
对 \(g\in C_b(Y)\)，复合函数 \(g\circ T\) 属于 \(C_b(X)\)。由推前积分公式，
\[
 \int_Y g\,\mathrm{d}(T_\#\mu_n)=\int_X g\circ T\,\mathrm{d}\mu_n
 \longrightarrow\int_X g\circ T\,\mathrm{d}\mu=\int_Y g\,\mathrm{d}(T_\#\mu).
\]
拓扑连续性也由同一等式给出：目标空间中每个测试积分与推前映射的复合，
都是源空间中的连续测试积分。
\end{proof}

称拓扑空间之间的连续映射 \(T:X\to Y\) 为\term[zhen-yingshe]{真映射}[proper map]，
若 \(Y\) 中每个紧集的原像都是 \(X\) 中的紧集。上面的连续映射原理不要求
\(T\) 具有这一性质，因为有界连续函数的复合仍有界连续；若改用 \(C_c(Y)\) 中的
紧支集测试，则 \(g\circ T\) 的支集一般未必紧。因而用淡收敛讨论推前测度时，
通常需要真映射或其他能够控制复合函数支集的条件；\chapref{b1:ch:11}的相应习题采用的正是这一条件。
在概率论中，若随机变量 \(Z_n\) 的分布为 \(\mu_n\)，它们分布的弱收敛也称为%
\term[yifenbu-shoulian]{依分布收敛}[convergence in distribution]，通常记作
\(Z_n\xrightarrow{d}Z\)，其中 \(Z\) 的分布为 \(\mu\)。
这一说法只比较分布，不要求随机变量处在同一个概率空间上。

\ProseParagraphBreak

后面的概率论版本允许 \(X\) 为无限维度量空间。在无限维 Hilbert 空间中，
连续紧支集函数可能只有零函数，\chapref{b1:ch:11}已说明其原因。
因此不能先用 \(C_c(X)\) 定义所有概率测度的收敛，再期待局部紧情形的论证自动推广。
而 \(C_b(X)\) 测试则在这一情形仍有充分的区分能力。

\ProseParagraphBreak

\cref{b1:tab:vague-weak-measures}概括两种收敛最直接的分界。淡收敛只测试局部可见的质量；弱收敛还用常函数 \(1\) 测试总质量，因此概率测度的弱极限不会丢失质量。
\begin{table}[htbp]
  \centering
  \caption[淡收敛与弱收敛]{淡收敛与弱收敛。表中淡收敛一行以局部紧 Hausdorff 空间上的 Radon 测度为背景。}
  \label{b1:tab:vague-weak-measures}
  \begin{tabularx}{\textwidth}{@{}p{2.4cm}p{4.0cm}X@{}}
    \toprule
    收敛方式 & 测试函数 & 对逃向无穷远质量的反应 \\
    \midrule
    淡收敛 & \(C_c(X)\) & 只观察固定紧区域，因而可能看不见逃逸质量 \\
    弱收敛 & \(C_b(X)\) & 常函数 \(1\) 同时控制总质量；概率极限仍是概率测度 \\
    \bottomrule
  \end{tabularx}
\end{table}

\cref{b1:fig:visual-f07}以原子位置说明测试类与紧族条件的作用。
% F07：本书原创矢量示意；依据所在节的定义、证明或明确模型。
\begin{figure}[htbp]
\centering
\begin{tikzpicture}[x=1cm,y=1cm,>=stealth,every node/.style={font=\small},line width=.65pt]

\draw[->,gray](0,0)--(5.7,0)node[right]{\(x\)};
\foreach \x/\n in {3/1,1.5/2,.75/4}{\draw[->,very thick,AnalysisBlue](\x,0)--(\x,1.8);\node[above]at(\x,1.8){\(1/\n\)};}
\draw[->,orange!80!black](2.8,2.6)--(.3,2.6);
\node at(2.5,3.25){\(\delta_{1/n}\Rightarrow\delta_0\)};
\draw[->,gray](7,0)--(12.7,0)node[right]{\(x\)};
\fill[AnalysisBlue!10](7,0)rectangle(9.3,1.4);
\node at(8.15,.65){固定紧区间};
\foreach \x in {8,10,12}{\draw[->,very thick,AnalysisBlue](\x,0)--(\x,1.8);}
\draw[->,orange!80!black](9,2.6)--(12.5,2.6);
\node at(9.8,3.25){\(\delta_n\) 淡收敛到零};
\node at(9.8,-.65){总质量始终为 \(1\)};

\end{tikzpicture}
\caption[原子测度的汇聚与质量逃逸]{原子测度的汇聚与质量逃逸。每根箭头都代表单位质量。逃逸族的紧支集测试趋零，但常函数一的积分不变，因此它不弱收敛到零测度。}
\label{b1:fig:visual-f07}
\end{figure}


% !TeX root = ../../Book1.tex
\section{弱收敛的等价刻画}
\label{b1:sec:1202}

弱收敛的定义要求逐个测试有界连续函数，但在具体问题中，集合的质量常常更容易计算。
从函数回到集合时，开集和闭集给出的不是同一个方向的不等式。
例如 \(\delta_{1/n}\Rightarrow\delta_0\)，而开集 \((0,1)\) 的质量从 \(1\) 降到 \(0\)，
闭集 \(\{0\}\) 的质量却从 \(0\) 升到 \(1\)。
这两种现象并不违背弱收敛；它们说明不能要求所有示性函数的积分都收敛。

本节固定度量空间 \((X,d)\)，并以概率测度为主要对象。
证明只使用距离函数、有限测度的收敛定理和集合的内外逼近，不预设空间局部紧或完备。
半连续积分与测度拓扑的关系可对照 Fremlin 的《Measure Theory》%
\cite[437J(f)--(g)、437K]{Fremlin2016Measure}；其中拓扑命名与本书的区别已在上一节说明。
% 实际阅读：Fremlin 官方 chap43.pdf 437J(c)--(g)、437K。
% 以下度量概率版本独立展开：距离逼近开闭集、分层积分与 Fatou、边界控制。
% Billingsley 仅作章首伴读建议，不声称本证明取自未阅读全文的版本。

\subsection{连续测试函数判据}
\label{b1:subsec:120201}

\begin{definition}\label{b1:def:measure-continuity-set}
设 \(\mu\) 为 \(X\) 上的 Borel 测度，\(A\) 为 Borel 集。
若 \(\mu(\partial A)=0\)，则称 \(A\) 为 \(\mu\) 的%
\term[cedu-lianxuji]{连续集}[continuity set]。
这里 \(\partial A=\overline A\setminus A^\circ\) 是相对于 \(X\) 的边界。
\end{definition}

定义中的连续不是说集合本身是一条曲线或一个连续像，而是说在这个集合的边界上，
极限测度没有质量。连续测试函数对示性函数的逼近误差便可以集中在越来越小的边界邻域中。

\begin{theorem}[Portmanteau]\label{b1:thm:portmanteau}
设 \(\mu_n,\mu\in\mathcal P(X)\)。下列条件等价：
\begin{equivalences}
\item\label{b1:item:portmanteau-continuous}
\(\mu_n\Rightarrow\mu\)，即所有有界连续函数的积分收敛。
\item\label{b1:item:portmanteau-open}
对每个开集 \(G\subset X\)，有 \(\mu(G)\leq\liminf_n\mu_n(G)\)。
\item\label{b1:item:portmanteau-closed}
对每个闭集 \(F\subset X\)，有 \(\limsup_n\mu_n(F)\leq\mu(F)\)。
\item\label{b1:item:portmanteau-continuity-set}
对每个 \(\mu\) 的连续集 \(A\)，有 \(\mu_n(A)\to\mu(A)\)。
\end{equivalences}
\end{theorem}
\begin{proof}
先设弱收敛成立。对非空闭集 \(F\)，令
\(h_k(x)=\max\{1-kd(x,F),0\}\)。这些函数属于 \(C_b(X)\)，并递减到 \(\mathbf1_F\)。
固定 \(k\)，由 \(\mathbf1_F\leq h_k\) 得
\[
 \limsup_{n\to\infty}\mu_n(F)
 \leq\lim_{n\to\infty}\int_X h_k\,d\mu_n
 =\int_X h_k\,d\mu.
\]
再令 \(k\to\infty\)，支配收敛定理给出闭集判据。空闭集另直接处理。
由于所有测度的总质量均为 \(1\)，对 \(G=X\setminus F\) 取补集可知，
开集判据与闭集判据等价。

现在只假设开集判据成立。若 \(f\in C_b(X)\) 且 \(0\leq f\leq1\)，
则 \(\{f>t\}\) 对每个 \(t\) 都开。由分层积分公式和 Fatou 引理，
\[
 \begin{aligned}
 \int_X f\,d\mu
 &=\int_0^1\mu(\{f>t\})\,dt\\
 &\leq\int_0^1\liminf_{n\to\infty}\mu_n(\{f>t\})\,dt
 \leq\liminf_{n\to\infty}\int_X f\,d\mu_n.
 \end{aligned}
\]
同样把这个结论用于 \(1-f\)，便得到反向的上极限不等式，故积分收敛。
任意实值有界连续函数经过平移和正数倍缩放可以落在 \([0,1]\) 中，
常函数的积分又都相同，因此得到弱收敛。

开闭集判据还给出：对任意 Borel 集 \(A\)，
\[
 \mu(A^\circ)\leq\liminf_n\mu_n(A)
 \leq\limsup_n\mu_n(A)\leq\mu(\overline A).
\]
若 \(\mu(\partial A)=0\)，左右两端都等于 \(\mu(A)\)，所以连续集判据成立。

最后假设连续集判据成立，取实值 \(f\in C_b(X)\)。集合
\(\{t:\mu(\{f=t\})>0\}\) 至多可数：对每个正整数 \(k\)，
其中质量不小于 \(1/k\) 的互不相交水平集至多有 \(k\) 个。
给定 \(\varepsilon>0\)，可选有限分点
\[
 t_0<\inf_X f\leq\sup_X f<t_m,
 \quad t_{j+1}-t_j<\varepsilon,
 \quad \mu(\{f=t_j\})=0.
\]
因为 \(f\) 连续，集合 \(A_j=\{t_j<f\leq t_{j+1}\}\) 的边界包含于
\(\{f=t_j\}\cup\{f=t_{j+1}\}\)，所以各 \(A_j\) 都是连续集。
令 \(s=\sum_{j=0}^{m-1}t_j\mathbf1_{A_j}\)，则 \(\|f-s\|_\infty<\varepsilon\)，且
\[
 \int_X s\,d\mu_n
 =\sum_{j=0}^{m-1}t_j\mu_n(A_j)
 \longrightarrow\sum_{j=0}^{m-1}t_j\mu(A_j)=\int_X s\,d\mu.
\]
两个概率测度下的一致逼近误差各小于 \(\varepsilon\)。先取积分差的上极限，
再令 \(\varepsilon\downarrow0\)，便得到 \(\int f\,d\mu_n\to\int f\,d\mu\)。
\end{proof}

这组等价判据通常称为\term[portmanteau-dingli]{Portmanteau 定理}[Portmanteau theorem]。
证明中两类逼近各有用途：距离函数从外逼近闭集，取值区间上的阶梯函数则逼近测试函数。
后一种逼近不是在空间 \(X\) 中划分成许多小区域，而是在函数的值域中分层。
因此即使 \(X\) 不紧，仍能用有限个水平集处理一个有界函数。

\subsection{开集判据}
\label{b1:subsec:120202}

开集判据只给出下界。若极限测度在一个开集内有质量，
那么这部分质量不能在近似测度中一直缺失；但近似测度中的额外质量可以靠近边界，
最终落到开集之外。为避免记反方向，可以先用 \(\delta_{1/n}\Rightarrow\delta_0\)
和开集 \((0,1)\) 检查一遍。

集合判据还有一个直接的函数版本。

\begin{proposition}\label{b1:prop:weak-lower-semicontinuity}
若 \(\mu_n\Rightarrow\mu\)，且 \(f:X\to[0,+\infty]\) 下半连续，则
\[
 \int_X f\,d\mu\leq\liminf_{n\to\infty}\int_X f\,d\mu_n.
\]
若 \(f\) 为有统一有限下界的实值或扩展实值下半连续函数，结论同样成立。
\end{proposition}
\begin{proof}
下半连续性正是说 \(\{f>t\}\) 对每个实数 \(t\) 都开。
对非负 \(f\)，重复前一定理中的分层积分与 Fatou 论证，只需把积分区间换为
\([0,+\infty)\)。即使某个积分为无穷，非负函数的这两个定理仍适用。
若 \(f\geq a> -\infty\)，先对 \(f-a\) 应用结论，再利用所有测度质量均为一，
把常数 \(a\) 加回即可。
\end{proof}

例如在 \(\mathbb R^d\) 上，\(x\mapsto|x|^p\) 对 \(p>0\) 连续且非负，故
\[
 \int_{\mathbb R^d}|x|^p\,d\mu
 \leq\liminf_n\int_{\mathbb R^d}|x|^p\,d\mu_n.
\]
因此统一的矩上界会传给弱极限，但矩本身未必收敛。
对 \(\mu_n=(1-1/n)\delta_0+(1/n)\delta_n\)，任意有界连续测试的积分趋于 \(f(0)\)，
而一阶绝对矩始终为 \(1\)，极限 \(\delta_0\) 的一阶绝对矩却为 \(0\)。
趋于零的质量仍可能带着不趋于零的矩，这与有界测试的定义相容。

后面谈弱紧集时，还需要拓扑形式，而不只是序列形式。对开集 \(G\ne X\)，令
\[
 g_k(x)=\min\{1,kd(x,X\setminus G)\}.
\]
有 \(g_k\uparrow\mathbf1_G\)，从而
\[
 \nu(G)=\sup_{k\geq1}\int_X g_k\,d\nu.
\]
右侧是弱连续函数的上确界，故 \(\nu\mapsto\nu(G)\) 在 \(\mathcal P(X)\) 上下半连续。
当 \(G=X\) 时该映射恒为 \(1\)，结论仍成立。特别地，
\(\{\nu:\nu(G)>a\}\) 是弱开集。这一表述能直接用于开覆盖与有限子覆盖的论证，
不必先假设测度空间的拓扑可由序列刻画。

\subsection{闭集判据}
\label{b1:subsec:120203}

若 \(F\) 闭，映射 \(\nu\mapsto\nu(F)\) 上半连续，
因为它等于 \(1-\nu(X\setminus F)\)。因此 \(\{\nu:\nu(F)\geq a\}\) 弱闭。
这说明“至少有给定质量留在一个闭集内”的条件能够传给极限。
相反，“质量恰好落在一个开集内”的条件不一定封闭。

\begin{proposition}\label{b1:prop:weak-closed-support}
设 \(F\subset X\) 闭。所有满足 \(\nu(F)=1\) 的概率测度构成 \(\mathcal P(X)\) 的弱闭子集。
若 \(\mu_n\Rightarrow\mu\)，且 \(\mu_n(F)=1\) 对每个 \(n\) 成立，则 \(\mu(F)=1\)。
\end{proposition}
\begin{proof}
由上一段的上半连续性，集合 \(\{\nu:\nu(F)\geq1\}\) 弱闭。
概率测度的值不超过一，所以该集合恰好是题述集合。序列结论是其直接后果。
\end{proof}

闭性在这里不能去掉：\(\delta_{1/n}\) 全部集中在开集 \((0,+\infty)\) 内，
弱极限却集中在其边界。也不能仅凭每一项分别集中在某个紧集内，就推出存在一个
对全列有效的紧集；紧集随 \(n\) 远移，仍会发生质量逃逸。

闭集方向也对应上半连续测试函数。
若 \(g:X\to[-\infty,+\infty)\) 上半连续且有有限上界 \(b\)，则
\(b-g\) 非负下半连续，前一命题给出
\[
 \limsup_{n\to\infty}\int_X g\,d\mu_n\leq\int_X g\,d\mu.
\]
这里允许积分为 \(-\infty\)，因为正部有界，不会出现未定义的
\(+\infty-\infty\)。若测试函数既无上界也无下界，则不能直接使用这两个不等式，
必须先解决可积性和尾部控制。

对于有限正测度，Portmanteau 判据需同时保留总质量收敛。
若 \(m_n=\mu_n(X)\to m=\mu(X)>0\)，把充分靠后的测度除以各自质量，
便归约为概率测度版本。若 \(m=0\)，则 \(\mu=0\)，且
\(\left|\int f\,d\mu_n\right|\leq\|f\|_\infty m_n\to0\)。
因此开集下界与总质量收敛合在一起，或闭集上界与总质量收敛合在一起，
都刻画有限正测度的弱收敛。单有开集下界或闭集上界则不够：
常值列 \(2\delta_0\) 与候选极限 \(\delta_0\) 满足前者，
常值列 \(\delta_0\) 与候选极限 \(2\delta_0\) 满足后者，但都不弱收敛。

\subsection{边界零测集判据}
\label{b1:subsec:120204}

连续集判据把开闭集的两个不等式合成了等式，但条件取决于极限测度，而不是近似测度。
在实轴上，\((a,b]\) 的边界是 \(\{a,b\}\)；
只要 \(\mu(\{a\})=\mu(\{b\})=0\)，就有
\(\mu_n\bigl((a,b]\bigr)\to\mu\bigl((a,b]\bigr)\)。
同样对边界为 \(\{b\}\) 的半直线 \(( -\infty,b]\) 使用判据，
便得到分布函数在极限分布连续点处的收敛。
这一条件反过来也能推出弱收敛，具体证明留在习题\ref{b1:ex:12-distribution-functions}中。

连续集还允许把测试函数适当放宽：函数可以不处处连续，只要不连续点不承载极限质量。

\begin{proposition}\label{b1:prop:weak-ae-continuous-test}
设 \(\mu_n\Rightarrow\mu\)，\(f:X\to\mathbb R\) 为有界 Borel 函数，
其不连续点集记为 \(D_f\)。若 \(\mu(D_f)=0\)，则
\[
 \int_X f\,d\mu_n\longrightarrow\int_X f\,d\mu.
\]
\end{proposition}
\begin{proof}
先说明零测条件的可测性。定义振幅
\[
 \omega_f(x)=\inf_{r>0}\sup\{\,|f(y)-f(z)|:y,z\in B(x,r)\,\}.
\]
若 \(\omega_f(x)<a\)，可选一个小球，使其中的振幅小于 \(a\)；
该球内充分靠近 \(x\) 的点也有同样性质。因此 \(\{\omega_f<a\}\) 开，
而 \(D_f=\bigcup_{k\geq1}\{\omega_f\geq1/k\}\) 是 Borel 集。

若 \(x\notin D_f\) 且 \(f(x)\ne t\)，连续性使 \(x\) 附近的函数值始终位于 \(t\) 的同一侧，
所以
\[
 \partial\{f\leq t\}\subseteq D_f\cup\{f=t\}.
\]
除去至多可数个具有正质量的水平值之后，\(\{f\leq t\}\) 便是 \(\mu\) 的连续集。
按定理\ref{b1:thm:portmanteau}证明中的方法，在 \(f\) 的有界值域两侧选端点，
并在内部选间隔小于 \(\varepsilon\) 的有限分点，使各分点都避开这些水平值。
每个集合 \(\{t_j<f\leq t_{j+1}\}\) 的边界包含于两端水平集与 \(D_f\) 的并，
因而同样是连续集。相应阶梯函数的积分收敛，且它对 \(f\) 的一致误差小于
\(\varepsilon\)。两次积分误差均小于 \(\varepsilon\)，令其趋于零即可。
\end{proof}

取 \(f=\mathbf1_A\)，其不连续点集正是 \(\partial A\)，便回到连续集判据。
取一个分段连续的有界函数，则只需检查分段边界上的极限质量。
例如 \(\mu_n\Rightarrow\mu\) 且 \(\mu(\{0\})=0\) 时，
\(\int\operatorname{sgn}(x)\,d\mu_n\to\int\operatorname{sgn}(x)\,d\mu\)。
若极限在零点有原子，这个结论可能失败。

至此我们已有定义、集合判据和半连续函数判据，但还没有保证任意给定的测度列存在弱极限。
下一节转向这个存在性问题所需的质量控制；在最后一节，紧族条件才会与测度空间中的紧性联系起来。

% !TeX root = ../../Book1.tex
\section{紧性}
\label{b1:sec:1203}

淡收敛只用紧支集函数观察测度。若质量逐渐移到每个固定紧集之外，这些测试函数就可能看不到它。要把局部测试推广到全部有界连续函数，需要控制被截断部分的质量，而且这种控制必须对整族测度一致。本节先定义这种条件，再说明它怎样补足淡收敛与弱收敛之间的差别。

\subsection{紧族}
\label{b1:subsec:120301}

\begin{definition}\label{b1:def:tight-family}
设 \(X\) 为 Hausdorff 空间，\(\mathcal F\) 是其上的一族有限正 Borel 测度。若对每个 \(\varepsilon>0\)，都存在紧集 \(K\subseteq X\)，使
\[
 \sup_{\mu\in\mathcal F}\mu(X\setminus K)<\varepsilon,
\]
则称 \(\mathcal F\) 为\term[jinzu]{紧族}[tight family]。%
\footnote{严加安的《测度论讲义》第二版\cite[第 6.3 节]{Yan2004Cedulun}使用“胎紧”\termidx[taijin]{胎紧}这一名称。本书在测度族的语境中固定使用“紧族”。}
若单元素族 \(\{\mu\}\) 满足此条件，则称测度 \(\mu\) 是紧的。
\end{definition}

定义中的量词顺序是先给误差，再为所有测度选同一个紧集；不能让这个紧集随 \(\mu\) 改变。每个有限 Radon 测度都是紧的，因为其总质量可以用紧集从内逼近。有限个这样的测度也组成紧族：对同一误差分别取紧集，再取它们的有限并。对无限族，这个论证便不再适用。

紧族条件不包含总质量一致有界。例如，在 \(\mathbb R\) 上，\(\{n\delta_0:n\geq1\}\) 的全部质量都在紧集 \(\{0\}\) 中，因此它是紧族，但总质量没有上界。概率测度的总质量固定为 \(1\)，才使这一额外问题自动消失。

“紧族”也不直接表示测度空间中的拓扑紧集。在 \(\mathbb R\) 上，族 \(\mathcal F=\{\delta_{1/n}:n\geq1\}\) 的支集都包含在紧集 \(\{0\}\cup\{1/n:n\geq1\}\) 中，因此是紧族。可是连续测试函数 \(f(x)=x/(1+|x|)\) 所给映射 \(\mu\mapsto\int f\,\mathrm d\mu\) 把 \(\mathcal F\) 映为 \(\{1/(n+1):n\geq1\}\)，后者不是 \(\mathbb R\) 中的紧集。故 \(\mathcal F\) 不是弱拓扑下的紧集。这里缺少的是极限 \(\delta_0\)；后面的相对紧性结论会涉及测度族的闭包，不能把闭包二字略去。

实际验证紧族时，常用一个函数惩罚远处的质量。该函数的低值区域必须真正紧，而不只是有界。

\begin{proposition}\label{b1:prop:tight-coercive-bound}
设 \(X\) 为 Hausdorff 空间，\(\Phi:X\rightarrow[0,\infty]\) 是 Borel 可测函数，且对每个 \(R>0\)，下水平集
\[
 K_R=\{\,x\in X\mid\Phi(x)\leq R\,\}
\]
都紧。若有限正 Borel 测度族 \(\mathcal F\) 满足
\[
 C=\sup_{\mu\in\mathcal F}\int_X\Phi\,\mathrm d\mu<\infty,
\]
则 \(\mathcal F\) 是紧族，并且 \(\sup_{\mu\in\mathcal F}\mu(X\setminus K_R)\leq C/R\)。
\end{proposition}
\begin{proof}
逐点不等式 \(R\mathbf1_{\{\Phi>R\}}\leq\Phi\) 给出
\[
 R\mu(X\setminus K_R)
 \leq\int_X\Phi\,\mathrm d\mu\leq C.
\]
对 \(\mu\) 取上确界，再选 \(R>C/\varepsilon\)，便得到定义中的紧集及严格小于 \(\varepsilon\) 的尾部界。若 \(C=0\)，任取 \(R>0\) 即可。
\end{proof}

在 \(\mathbb R^d\) 上，取 \(\Phi(x)=|x|^p\)，其中 \(p>0\)。其下水平集是闭球，故任何满足一致 \(p\) 阶矩界 \(\int|x|^p\,\mathrm d\mu\leq C\) 的测度族都是紧族，具体估计为
\[
 \mu\bigl(\{\,x:|x|>r\,\}\bigr)\leq Cr^{-p}
 \quad(r>0).
\]
这给出了可计算的紧集半径。不过，矩界是充分条件，不是紧性的必要条件。例如，实线上密度为 \(1/\bigl(\pi(1+x^2)\bigr)\) 的概率测度满足
\[
 \mu\bigl(\mathbb R\setminus[-R,R]\bigr)
 \leq\frac{2}{\pi R}
 \quad(R>0),
\]
因为在 \(x>R\) 上 \((1+x^2)^{-1}\leq x^{-2}\)。所以它是紧的，而一阶绝对矩发散：对 \(x\geq1\)，有 \(x/(1+x^2)\geq1/(2x)\)。

紧下水平集这一假设也不能删去。在 \(\ell^2\) 中取标准正交向量 \(e_n\) 及点质量 \(\delta_{e_n}\)。虽然 \(\int\|x\|^2\,\mathrm d\delta_{e_n}=1\)，这些测度不组成紧族。任一紧集至多包含有限个 \(e_n\)：否则，由 \(\|e_n-e_m\|=\sqrt2\) 可知其中不存在 Cauchy 子列，与紧度量空间的序列紧性矛盾。于是每个紧集之外都有某个 \(\delta_{e_n}\) 的全部质量。有限维中可用的“矩控制尾部”仍成立，但其尾部是球外部分；无限维闭球一般不紧。

\subsection{质量逃逸}
\label{b1:subsec:120302}

质量逃逸既可能来自支集整体平移，也可能来自质量不断摊薄。后一种情形中，支集甚至始终包含原点。令实线上的概率测度
\[
 \lambda_n=\frac1{2n}\mathbf1_{[-n,n]}m,
\]
其中 \(m\) 是 Lebesgue 测度。给定 \(f\in C_c(\mathbb R)\)，取 \(R>0\) 使 \(\operatorname{supp}f\subseteq[-R,R]\)。当 \(n\geq R\) 时，
\[
 \left|\int_{\mathbb R}f\,\mathrm d\lambda_n\right|
 \leq\frac{R}{n}\|f\|_\infty\longrightarrow0.
\]
因此 \(\lambda_n\xrightarrow{v}0\)，但不可能弱收敛到零测度，因为常函数 \(1\) 的积分始终为 \(1\)。对任何紧集 \(K\subseteq[-R,R]\)，还有
\[
 \lambda_n(\mathbb R\setminus K)\geq1-\frac{R}{n}
 \quad(n\geq R),
\]
故紧族条件失败。每个固定观察区域中的质量都趋于零，而总质量没有消失，只是不再停留在任何统一的紧区域内。

这一例子也可以只让一部分质量逃逸。固定 \(0<\alpha<1\)，令
\[
 \mu_n=(1-\alpha)\delta_0+\alpha\lambda_n.
\]
对每个紧支集测试函数直接积分，得到 \(\mu_n\xrightarrow{v}(1-\alpha)\delta_0\)。所有 \(\mu_n\) 都是概率测度，极限却只有质量 \(1-\alpha\)。因而问题不只是“极限是不是零”，而是极限是否保留了原来的全部质量。

反过来，支集无界增长本身不妨碍紧族条件。令
\[
 \eta_n=\left(1-\frac1n\right)\delta_0+\frac1n\delta_n.
\]
给定 \(\varepsilon>0\)，选 \(N\) 使 \(1/N<\varepsilon\)，并取紧集 \(K=\{0,1,\ldots,N\}\)。当 \(n\leq N\) 时尾部质量为零，当 \(n>N\) 时尾部质量为 \(1/n<\varepsilon\)。因此 \(\{\eta_n:n\geq1\}\) 是紧族。实际上，对任意 \(f\in C_b(\mathbb R)\)，
\[
 \left|\int f\,\mathrm d\eta_n-f(0)\right|
 =\frac1n\bigl|f(n)-f(0)\bigr|
 \leq\frac2n\|f\|_\infty,
\]
所以 \(\eta_n\Rightarrow\delta_0\)。可以远去的是支集中的点；紧族条件约束的是这些点所携带的质量。

\subsection{淡收敛与弱收敛的差别}
\label{b1:subsec:120303}

以下回到局部紧 Hausdorff 空间，并采用\cref{b1:def:vague-convergence}的 \(C_c(X)\) 测试和\cref{b1:def:measure-weak-convergence}的 \(C_b(X)\) 测试。有限性与正性在本小节中各有作用：前者使常数函数可积，后者使总质量减去局部质量能够控制尾部。

% 实读来源：Fremlin, Measure Theory, vol.4，官方 chap43.pdf，
% 437O--P，印刷 pp.66--67：一致紧性定义及其与相对紧性的关系；
% 437J(c)--(d)，印刷 p.61：vague/narrow 拓扑的原书约定。
% https://www1.essex.ac.uk/maths/people/fremlin/chap43.pdf
% Fremlin的vague以Cb测试，不等于本书以Cc测试的淡收敛。
% 下述LCH序列比较直接用第11章截断引理证明，不借437P或后节Prokhorov。
% 能量判据、摊薄/部分逃逸及符号相消例由定义和积分估计自行组织。
测度族的一致紧性及其与拓扑紧性的关系，可参读 D.~H. Fremlin 的《Measure Theory》\cite[第 437O--P 节]{Fremlin2016Measure}。比较该书与本章时须留意，它在第 437J 节把由 \(C_b(X)\) 测试生成的拓扑称为 vague topology，并非本书的淡收敛约定。下面的结论直接用连续截断证明。

\begin{theorem}\label{b1:thm:vague-tight-weak}
设 \(X\) 为局部紧 Hausdorff 空间，\(\mu_n,\mu\) 均为有限正 Radon 测度，且 \(\mu_n\xrightarrow{v}\mu\)。则下列条件等价：
\begin{equivalences}
\item\label{b1:item:vague-tight-family}\(\{\mu_n:n\geq1\}\) 是紧族；
\item\label{b1:item:vague-total-mass}\(\mu_n(X)\rightarrow\mu(X)\)；
\item\label{b1:item:vague-to-weak}\(\mu_n\Rightarrow\mu\)。
\end{equivalences}
\end{theorem}
\begin{proof}
先由\itemref{b1:item:vague-tight-family}推出\itemref{b1:item:vague-to-weak}。固定 \(f\in C_b(X)\) 及 \(\varepsilon>0\)。由紧族条件及 \(\mu\) 的有限 Radon 性，可以选同一个紧集 \(K\)，使
\[
 \sup_n\mu_n(X\setminus K)<\varepsilon,
 \quad \mu(X\setminus K)<\varepsilon.
\]
由\cref{b1:lem:lch-cutoff}取 \(h\in C_c(X)\)，满足 \(0\leq h\leq1\) 且 \(h=1\) 于 \(K\)。函数 \(fh\) 有紧支集，故其积分由淡收敛控制；其余部分满足
\begin{align*}
 \left|\int_X f\,\mathrm d\mu_n-\int_X f\,\mathrm d\mu\right|
 &\leq\left|\int_X fh\,\mathrm d\mu_n-\int_X fh\,\mathrm d\mu\right|\\
 &\quad+\|f\|_\infty\bigl(\mu_n(X\setminus K)+\mu(X\setminus K)\bigr).
\end{align*}
先令 \(n\rightarrow\infty\)，再令 \(\varepsilon\downarrow0\)，右端趋于零，得到弱收敛。由\itemref{b1:item:vague-to-weak}推出\itemref{b1:item:vague-total-mass}只需取测试函数 \(f=1\)。

最后假设\itemref{b1:item:vague-total-mass}成立。给定 \(\varepsilon>0\)，先取紧集 \(K\)，使 \(\mu(X\setminus K)<\varepsilon/4\)，再取上述截断 \(h\)，令 \(L=\operatorname{supp}h\)。由 \(h=1\) 于 \(K\) 及 \(0\leq h\leq1\)，
\[
 0\leq\mu(X)-\int_Xh\,\mathrm d\mu<\frac{\varepsilon}{4}.
\]
总质量收敛与淡收敛一起给出
\[
 \mu_n(X)-\int_Xh\,\mathrm d\mu_n
 \longrightarrow\mu(X)-\int_Xh\,\mathrm d\mu.
\]
由于 \(h=0\) 于 \(X\setminus L\)，存在 \(N\)，使所有 \(n\geq N\) 都满足
\[
 \mu_n(X\setminus L)
 \leq\int_X(1-h)\,\mathrm d\mu_n<\frac{\varepsilon}{2}.
\]
对有限个首项 \(\mu_1,\ldots,\mu_{N-1}\)，分别用有限 Radon 性取紧集 \(K_j\)，使 \(\mu_j(X\setminus K_j)<\varepsilon/2\)。将这些紧集与 \(L\) 合并为 \(K'\)，便有 \(\sup_n\mu_n(X\setminus K')\leq\varepsilon/2<\varepsilon\)。因此整列组成紧族。
\end{proof}

证明中真正把远处质量收回来的量是 \(\mu_n(X)-\int h\,\mathrm d\mu_n\)。它同时使用常数函数和紧支集函数，因而不能由淡收敛单独控制。最后补入有限个首项则用到了序列的结构；对一般网，不能把某个指标之前的所有项当作有限集，这也是将拓扑紧性与序列论证分开的一个理由。

\begin{corollary}\label{b1:cor:vague-probability-weak}
在局部紧 Hausdorff 空间上，若 Radon 概率测度列淡收敛到一个 Radon 概率测度，则它也弱收敛到该测度，并且整列组成紧族。
\end{corollary}
\begin{proof}
所有测度的总质量均为 \(1\)，故前一定理的总质量条件自动成立。
\end{proof}

必须要求淡极限仍是概率测度；本节的 \(\lambda_n\) 和 \(\mu_n\) 已分别给出全部与部分质量丢失的情形。若 \(X\) 本身紧，则 \(C_c(X)=C_b(X)=C(X)\)，这两种收敛从定义上就相同；非紧空间中，紧族条件起到了统一限制尾部的作用。

对符号测度和复测度，尾部应由全变差衡量，而不是用可能发生抵消的测度值衡量。

\begin{proposition}\label{b1:prop:signed-tight-vague-weak}
设 \(X\) 为局部紧 Hausdorff 空间，\(\nu_n,\nu\) 是有限符号或复 Radon 测度。若 \(\nu_n\xrightarrow{v}\nu\)，且正测度族 \(\{|\nu_n|:n\geq1\}\) 是紧族，则对每个 \(f\in C_b(X)\)，有
\[
 \int_Xf\,\mathrm d\nu_n\longrightarrow\int_Xf\,\mathrm d\nu.
\]
\end{proposition}
\begin{proof}
给定 \(\varepsilon>0\)，由假设及 \(|\nu|\) 的有限 Radon 性，取紧集 \(K\)，使 \(\sup_n|\nu_n|(X\setminus K)<\varepsilon\) 且 \(|\nu|(X\setminus K)<\varepsilon\)。再取 \(h\in C_c(X)\)，使 \(0\leq h\leq1\)、\(h=1\) 于 \(K\)。全变差的积分估计给出
\[
 \left|\int_X f(1-h)\,\mathrm d\nu_n\right|
 \leq\|f\|_\infty\cdot|\nu_n|(X\setminus K)
 <\varepsilon\|f\|_\infty,
\]
对 \(\nu\) 也有同样的估计，而 \(fh\) 的积分由淡收敛得到极限。分解 \(f=fh+f(1-h)\)，先令 \(n\rightarrow\infty\)，再令 \(\varepsilon\downarrow0\)，即得结论。
\end{proof}

单纯的总质量收敛在这里不够。在实线上令 \(\nu_n=\delta_n-\delta_{n+1/2}\)。对每个紧支集函数，其积分最终为零，因此 \(\nu_n\xrightarrow{v}0\)；同时 \(\nu_n(\mathbb R)=0\)，甚至 \(|\nu_n|(\mathbb R)=2\) 一致有界。然而取 \(f(x)=\cos(2\pi x)\)，便有 \(\int f\,\mathrm d\nu_n=2\)，并不趋于零。对任意固定紧集 \(K\)，当 \(n\) 足够大时，\(\nu_n(\mathbb R\setminus K)=0\)，但 \(|\nu_n|(\mathbb R\setminus K)=2\)。总质量的相消掩盖了正负两部分同时逃逸，只有全变差尾部能给出所需的误差控制。

% !TeX root = ../../Book1.tex
\section{Prokhorov 定理}
\label{b1:sec:1204}

% TODO: 撰写本节正文。

\subsection{Polish 空间}
\label{b1:subsec:120401}

待填充。

\subsection{紧性与相对紧性}
\label{b1:subsec:120402}

待填充。

\subsection{概率测度中的应用}
\label{b1:subsec:120403}

待填充。

\subsection{Radon 空间版本}
\label{b1:subsec:120404}

待填充。


\begin{chaptersummary}
测度收敛首先取决于测试函数。局部紧空间上的淡收敛只使用连续紧支集函数；
总质量一致有界时，它与 \(C_0(X)^*\) 中的弱-*收敛相接。
由 \(C_b(X)\) 测试的弱收敛还保留总质量，却不要求每个 Borel 集的质量逐一收敛。
Portmanteau 定理给出开集的下极限和闭集的上极限不等式，
而边界对极限测度为零时，两边才合成集合质量的收敛。

紧族条件使同一个紧集之外的质量对所有测度都小。
在有限正 Radon 测度已经淡收敛的前提下，这一条件、总质量收敛与弱收敛等价。
正性使总质量减去局部质量能够估计尾部；允许符号或复数质量以后，
相消会破坏这一估计，必须改由全变差控制。
对于无界测试，即使测度已经弱收敛，也仍需另查积分尾部，不能只检查质量尾部。

在波兰空间上，Prokhorov 定理把紧族、相对弱紧和收敛子列联系起来。
证明先在紧空间中取得极限，再用共同紧集说明极限仍集中在原空间。
紧性本身不决定极限是哪一个测度；实际应用中还要用分布函数、
可分离测试或不变性等条件辨认它。
一般完全正则空间上的 Radon 版本保留了紧族推出相对弱紧的方向，
反向涉及 Prokhorov 性质，而序列版本还需要额外的度量化条件。

至此，第二编已从赋范函数空间推进到测度构成的空间。
下一编转向小尺度的质量分布，覆盖工具将承担新的任务；
本章关于测试、截断与紧性的区分，仍可用于分析那些随尺度变化的测度。
\end{chaptersummary}

\BookExercises

% 八题均为自拟，逐题保留工具与改编目的说明；不冒称原书题号。

\begin{exercise}\label{b1:ex:12-distribution-functions}
设 \(\mu_n,\mu\) 为 \(\mathbb R\) 上的概率测度，记
\[
 F_n(t)=\mu_n\bigl((-\infty,t]\bigr),\quad
 F(t)=\mu\bigl((-\infty,t]\bigr).
\]
\begin{enumerate}
\item\label{b1:item:12-cdf-criterion}
证明 \(\mu_n\Rightarrow\mu\) 当且仅当 \(F_n(t)\rightarrow F(t)\) 对 \(F\) 的每个连续点 \(t\) 成立。
\item\label{b1:item:12-cdf-exponential}
设 \(\mu_n\) 关于 Lebesgue 测度的密度为
\[
 p_n(x)=\begin{cases}(1-x/n)^{n-1},&0<x<n,\\0,&\text{其余情形}.\end{cases}
\]
验证其总质量为 \(1\)，并求 \(\mu_n\) 的弱极限。
\end{enumerate}
\end{exercise}
% 来源：自拟；把Portmanteau 定理与实轴区间划分结合，并计算一个密度极限。
% 充分性实际证明有界连续测试的收敛，不把区间生成 Borel 集误作弱收敛判据。
\begin{solution}
分布函数 \(F\) 右连续，且 \(F(t)-F(t-)=\mu(\{t\})\)。所以 \(F\) 的连续点恰好是 \(\mu(\{t\})=0\) 的点。若 \(\mu_n\Rightarrow\mu\)，则集合 \((-\infty,t]\) 的边界为 \(\{t\}\)，\cref{b1:thm:portmanteau}的边界零测集判据给出必要性。

反过来，设连续点处的收敛成立。正质量的单点至多可数，故这些连续点稠密。给定 \(f\in C_b(\mathbb R)\) 和 \(\varepsilon>0\)，选连续点 \(a<b\)，使 \(F(a)<\varepsilon\)、\(1-F(b)<\varepsilon\)。则对充分大的 \(n\)，
\[
 \mu_n\bigl(\mathbb R\setminus(a,b]\bigr)<4\varepsilon,
 \quad\mu\bigl(\mathbb R\setminus(a,b]\bigr)<2\varepsilon.
\]
由 \(f\) 在 \([a,b]\) 上的一致连续性，可选全部由 \(F\) 连续点组成的划分
\(a=t_0<t_1<\cdots<t_k=b\)，使每个小区间上的振幅小于 \(\varepsilon\)。在 \((t_{j-1},t_j]\) 上用常数 \(f(t_j)\) 逼近 \(f\)。这些区间的质量满足
\[
 \mu_n\bigl((t_{j-1},t_j]\bigr)
 =F_n(t_j)-F_n(t_{j-1})
 \longrightarrow F(t_j)-F(t_{j-1}).
\]
有限个区间的阶梯函数积分因而收敛；两次区间内逼近与区间外尾部给出
\[
 \limsup_n\left|\int f\,\mathrm d\mu_n-\int f\,\mathrm d\mu\right|
 \le2\varepsilon+6\|f\|_\infty\varepsilon.
\]
令 \(\varepsilon\downarrow0\)，便得到弱收敛。

第二问中，直接积分得 \(\int_0^n(1-x/n)^{n-1}\,\mathrm dx=1\)。分布函数为
\[
 F_n(t)=\begin{cases}
 0,&t\le0,\\
 1-(1-t/n)^n,&0<t<n,\\
 1,&t\ge n.
 \end{cases}
\]
固定 \(t>0\)，有 \((1-t/n)^n\rightarrow e^{-t}\)，故极限分布函数为 \(F(t)=0\) 于 \(t\le0\)，\(F(t)=1-e^{-t}\) 于 \(t>0\)。它处处连续，对应密度 \(e^{-x}\mathbf1_{(0,\infty)}(x)\)。第一问给出所求弱极限。
\end{solution}

\begin{exercise}\label{b1:ex:12-unbounded-tests}
设 \(S\) 为波兰空间，\(\mu_n\Rightarrow\mu\) 为概率测度的弱收敛，\(f:S\rightarrow\mathbb R\) 连续但不必有界。假设
\[
 A(R):=\sup_n\int_{\{|f|>R\}}|f|\,\mathrm d\mu_n
 \longrightarrow0\quad(R\longrightarrow\infty).
\]
证明 \(f\) 关于 \(\mu\) 及每个 \(\mu_n\) 都可积，且
\[
 \int_Sf\,\mathrm d\mu_n\longrightarrow\int_Sf\,\mathrm d\mu.
\]
给出各测度都具有有限一阶矩的例子，说明若删除上述尾部条件，结论可能失败。
\end{exercise}
% 来源：自拟；将第五章幅度尾部一致可积的思想用于随 n 改变的测度。
% 解答把极限测度的尾部控制也证出，不只对近似列作截断。
\begin{solution}
选 \(R_0\) 使 \(A(R_0)<\infty\)，则
\(\int|f|\,\mathrm d\mu_n\le R_0+A(R_0)\)，故每项可积。令
\[
 f_R=\max\bigl\{-R,\min\{R,f\}\bigr\}.
\]
它有界连续，且 \(|f-f_R|=(|f|-R)_+\)。对固定的 \(R,M>0\)，函数 \(\min\{|f-f_R|,M\}\) 有界连续，因而
\[
 \int_S\min\{|f-f_R|,M\}\,\mathrm d\mu
 =\lim_n\int_S\min\{|f-f_R|,M\}\,\mathrm d\mu_n
 \le A(R).
\]
令 \(M\uparrow\infty\)，得到 \(\int|f-f_R|\,\mathrm d\mu\le A(R)\)。取 \(R=R_0\) 可知 \(f\in L^1(\mu)\)，再对充分大的 \(R\) 写
\[
 \left|\int f\,\mathrm d\mu_n-\int f\,\mathrm d\mu\right|
 \le2A(R)+\left|\int f_R\,\mathrm d\mu_n-\int f_R\,\mathrm d\mu\right|.
\]
先令 \(n\to\infty\)，再令 \(R\to\infty\)，结论即得。

令 \(U_a\) 表示区间 \([a,a+1]\) 上的均匀概率测度，并取
\[
 \mu_n=(1-1/n)U_0+(1/n)U_n.
\]
对 \(g\in C_b(\mathbb R)\)，积分与 \(U_0\) 的积分之差不超过 \(2\|g\|_\infty/n\)，故 \(\mu_n\Rightarrow U_0\)。然而对于 \(f(x)=x\)，
\[
 \int x\,\mathrm d\mu_n
 =(1-1/n)\frac12+\frac1n\left(n+\frac12\right)=\frac32,
 \quad\int x\,\mathrm dU_0=\frac12.
\]
对任意 \(R\)，取 \(n>R\)，远处区间对 \(A(R)\) 的贡献至少为 \(1\)，所以尾部条件确实不成立。
\end{solution}

\begin{exercise}\label{b1:ex:12-atomic-approximation}
设 \(S\) 为非空可分度量空间，\(D\subset S\) 是可数稠密集，\(\mu\) 为 Borel 概率测度。证明存在支集包含于 \(D\) 的有限原子概率测度 \(\mu_n\)，使 \(\mu_n\Rightarrow\mu\)。进一步证明，原子质量均可取为非负有理数。由此说明：即使给定了一个可数稠密点集，仍能只用该点集上的有限原子测度作弱逼近。
\end{exercise}
% 来源：自拟；用可数球覆盖、可数可加性与控制收敛实际构造逼近。
% 本题只需可分，不需要完备或局部紧；不把有限原子逼近误说成全变差逼近。
\begin{solution}
把 \(D\) 列为 \((d_j)\)，有限集的情形允许重复列点。对每个 \(n\)，半径 \(1/n\) 的这些开球覆盖 \(S\)。由测度对递增并的连续性，存在 \(N_n\)，使
\[
 E_n:=S\setminus\bigcup_{j=1}^{N_n}B(d_j,1/n)
 \quad\text{满足}\quad\mu(E_n)<2^{-n}.
\]
对 \(x\notin E_n\)，令 \(T_n(x)\) 为首个满足 \(d(x,d_j)<1/n\) 的 \(d_j\)；对 \(x\in E_n\)，令 \(T_n(x)=d_1\)。按首个序号分配得到 Borel 划分，故 \(T_n\) 为有限值 Borel 映射。测度 \(\mu_n=(T_n)_\#\mu\) 是所需形式的有限原子概率测度。

由于
\[
 \mu\left(\bigcup_{n\ge m}E_n\right)\le\sum_{n\ge m}2^{-n}\longrightarrow0,
\]
几乎每个点只属于有限多个 \(E_n\)。故对这些点，\(d\bigl(T_n(x),x\bigr)\rightarrow0\)。对每个 \(f\in C_b(S)\)，由连续性及控制收敛，
\[
 \int_Sf\,\mathrm d\mu_n
 =\int_Sf\circ T_n\,\mathrm d\mu
 \longrightarrow\int_Sf\,\mathrm d\mu.
\]

最后，设 \(\mu_n=\sum_{j=1}^ka_j\delta_{x_j}\)。对充分大的整数 \(Q\)，把前 \(k-1\) 个系数换成 \(b_j=\lfloor Qa_j\rfloor/Q\)，并令 \(b_k=1-\sum_{j<k}b_j\)。则各 \(b_j\ge0\)，总和为 \(1\)，且
\[
 \sum_{j=1}^k|a_j-b_j|\le\frac{2(k-1)}Q.
\]
选 \(Q\) 使右边小于 \(1/n\)，得到有理质量测度 \(\widetilde\mu_n\)。它与 \(\mu_n\) 对每个有界测试 \(f\) 的积分之差小于 \(\|f\|_\infty/n\)，所以 \(\widetilde\mu_n\Rightarrow\mu\)。全部这样的有理有限原子测度仅有可数多个。
\end{solution}

\begin{exercise}\label{b1:ex:12-discontinuous-mapping}
设 \(S,T\) 为波兰空间，\(\mu_n\Rightarrow\mu\)，且 \(F:S\rightarrow T\) 为 Borel 映射。记 \(D_F\) 为 \(F\) 的不连续点集。证明：若 \(\mu(D_F)=0\)，则 \(F_\#\mu_n\Rightarrow F_\#\mu\)。

再令 \(\mu_n\) 为 \([-1/n,1/n]\) 上的均匀概率测度，\(F=\mathbf1_{\{0\}}:\mathbb R\rightarrow\mathbb R\)。求 \(\mu_n\)、\(F_\#\mu_n\) 各自的弱极限，并比较后一极限与 \(F_\#\delta_0\)。
\end{exercise}
% 来源：自拟；从闭集版Portmanteau 定理证明几乎处处连续的映射版本，
% 新增不连续点集的处理；并非重复正文处处连续的推前结论。
\begin{solution}
度量空间映射的连续点集是可数个开集的交：对每个正整数 \(k\)，取所有使 \(F\) 的像直径小于 \(1/k\) 的开集之并，再对 \(k\) 取交即可。因此 \(D_F\) 是 Borel 集。

给定闭集 \(B\subset T\)，令 \(A=F^{-1}(B)\)。若 \(x\notin D_F\) 且 \(x\notin A\)，则 \(F(x)\notin B\)，连续性使 \(x\) 有一个与 \(A\) 不交的邻域。因此
\[
 \overline A\subset A\cup D_F.
\]
由\cref{b1:thm:portmanteau}的闭集不等式，
\[
 \limsup_n(F_\#\mu_n)(B)
 =\limsup_n\mu_n(A)
 \le\mu(\overline A)\le\mu(A)=(F_\#\mu)(B).
\]
各推前都是波兰空间上的 Borel 概率测度，故同一定理的闭集判据给出结论。

例子中，对 \(g\in C_b(\mathbb R)\)，
\[
 \left|\int g\,\mathrm d\mu_n-g(0)\right|
 \le\sup_{|x|\le1/n}|g(x)-g(0)|\longrightarrow0,
\]
所以 \(\mu_n\Rightarrow\delta_0\)。每个 \(\mu_n\) 都不给单点质量，故 \(F_\#\mu_n=\delta_0\)，而 \(F_\#\delta_0=\delta_1\)。映射在极限测度承载全部质量的点不连续，因而推前与极限不能交换。
\end{solution}

\begin{exercise}\label{b1:ex:12-product-limits}
设 \(S,T\) 为波兰空间，\(\mu_n\Rightarrow\mu\) 为 \(S\) 上概率测度的弱收敛，\(\nu_n\Rightarrow\nu\) 为 \(T\) 上概率测度的弱收敛。证明
\[
 \mu_n\otimes\nu_n\Rightarrow\mu\otimes\nu
\]
在 \(S\times T\) 上成立。证明应解释：为什么只算出 \(f(x)g(y)\) 这类可分离测试的积分还不够，以及紧性如何把这些测试提升为所有有界连续测试。
\end{exercise}
% 来源：自拟；用 Prokhorov 抽取聚点，再由闭矩形辨认唯一极限。
% 不预借 Stone--Weierstrass，不直接把可分离测试当成全部 Cb。
\begin{solution}
弱收敛列连同其极限构成弱紧集，故由\cref{b1:thm:prokhorov}，\(\{\mu_n\}\) 与 \(\{\nu_n\}\) 分别为紧族。对给定 \(\varepsilon>0\)，可取紧集 \(K\subset S\)、\(L\subset T\)，使两族在各自补集上的质量都小于 \(\varepsilon/2\)。于是
\[
 (\mu_n\otimes\nu_n)\bigl((S\times T)\setminus(K\times L)\bigr)
 \le\mu_n(S\setminus K)+\nu_n(T\setminus L)<\varepsilon.
\]
乘积测度也构成紧族。波兰空间之积仍为波兰空间，且其 Borel \(\sigma\)-代数等于乘积 Borel \(\sigma\)-代数，所以第\ref{b1:ch:07}章的乘积测度及 Prokhorov 定理都适用。

从任意子列抽出弱收敛子列，记极限为 \(\lambda\)。对 \(f\in C_b(S)\)、\(g\in C_b(T)\)，Fubini 定理给出
\[
 \int_{S\times T}f(x)g(y)\,\mathrm d\lambda
 =\left(\int_Sf\,\mathrm d\mu\right)
  \left(\int_Tg\,\mathrm d\nu\right).
\]
给定非空闭集 \(A\subset S\)、\(B\subset T\)，取
\[
 f_k(x)=\max\{1-kd(x,A),0\},\quad
 g_k(y)=\max\{1-kd(y,B),0\}.
\]
这些函数分别递减到 \(\mathbf1_A\)、\(\mathbf1_B\)。在上式中令 \(k\to\infty\)，由控制收敛得到
\(\lambda(A\times B)=\mu(A)\nu(B)\)；空集的情形直接成立。闭矩形构成生成乘积 Borel \(\sigma\)-代数的 \(\pi\) 系统，故有限测度的生成类唯一性给出 \(\lambda=\mu\otimes\nu\)。

现在若某个有界连续测试的积分不收敛到目标值，就能选出一条子列，使偏差始终大于某个正数。紧性又为它提供弱收敛子列，而刚才的辨认说明其极限只能是 \(\mu\otimes\nu\)，矛盾。因此全列弱收敛。可分离测试在证明中承担的是辨认聚点，而聚点的存在来自紧性。
\end{solution}

\begin{exercise}\label{b1:ex:12-moment-tightness}
设 \(\mathcal P\) 为 \(\mathbb R^d\) 上的概率测度族，且对某个 \(p>0\)，
\[
 \sup_{\rho\in\mathcal P}\int_{\mathbb R^d}\|x\|^p\,\mathrm d\rho(x)\le C<\infty.
\]
\begin{enumerate}
\item\label{b1:item:12-moment-tight}
证明 \(\mathcal P\) 为紧族，因而相对弱紧。
\item\label{b1:item:12-moment-limit}
若 \(\mu_n\in\mathcal P\) 且 \(\mu_n\Rightarrow\mu\)，证明 \(\int\|x\|^p\,\mathrm d\mu\le C\)，并对每个 \(0<q<p\) 证明 \(q\) 阶矩收敛。
\item\label{b1:item:12-moment-converse}
在 \(\mathbb R\) 上构造紧概率测度族，使其中每个测度的一阶矩都有限，但这些一阶矩没有共同上界。
\end{enumerate}
\end{exercise}
% 来源：自拟；矩估计用于紧性，较高阶矩用于幅度一致可积，两种用途分开。
% 逆向反例采用截断 Cauchy 密度，不重复上一题的少量质量远移模型。
\begin{solution}
由积分的基本估计，
\[
 \rho\{\,\|x\|>R\,\}\le\frac C{R^p}\quad(\rho\in\mathcal P).
\]
欧氏闭球紧，故 \(\mathcal P\) 满足\cref{b1:def:tight-family}，相对弱紧性由\cref{b1:thm:prokhorov}得到。

对任意 \(M>0\)，\(x\mapsto\min\{\|x\|^p,M\}\) 有界连续，所以其关于 \(\mu\) 的积分不超过 \(C\)。令 \(M\uparrow\infty\)，单调收敛给出极限的 \(p\) 阶矩上界。若 \(0<q<p\)，则
\[
 \int_{\{\|x\|^q>R\}}\|x\|^q\,\mathrm d\mu_n
 \le R^{-(p-q)/q}\int\|x\|^p\,\mathrm d\mu_n
 \le C R^{-(p-q)/q}.
\]
右端一致趋于零。将\cref{b1:ex:12-unbounded-tests}用于 \(f(x)=\|x\|^q\)，便得
\(\int\|x\|^q\,\mathrm d\mu_n\rightarrow\int\|x\|^q\,\mathrm d\mu\)。

对第三问，令 \(\rho_n\) 的密度为
\[
 r_n(x)=\frac{\mathbf1_{[-n,n]}(x)}{2\arctan n\,(1+x^2)}\quad(n\ge1).
\]
分母中的常数将积分归一化为 \(1\)。因 \(2\arctan n\ge\pi/2\)，对 \(R>0\)，
\[
 \sup_n\rho_n\{\,|x|>R\,\}
 \le\frac2\pi\int_{|x|>R}\frac{\mathrm dx}{1+x^2}
 \le\frac4{\pi R}.
\]
故 \(\{\rho_n\}\) 为紧族。然而
\[
 \int|x|\,\mathrm d\rho_n
 =\frac{\log(1+n^2)}{2\arctan n}\longrightarrow\infty.
\]
每项因支集有界而具有有限一阶矩，但紧性不控制该无界测试的积分。
\end{solution}

\begin{exercise}\label{b1:ex:12-holder-paths}
在实 Banach 空间 \(E=C([0,1];\mathbb R)\) 上取一致范数。固定 \(0<\alpha\le1\)，并令
\[
 H_\alpha(f)=\sup_{0\le s<t\le1}\frac{|f(t)-f(s)|}{|t-s|^\alpha}\in[0,\infty].
\]
\begin{enumerate}
\item\label{b1:item:12-holder-compact}
证明对每个 \(M>0\)，集合
\[
 K_M=\{\,f\in E\mid |f(0)|\le M,\ H_\alpha(f)\le M\,\}
\]
紧。请直接给出有限网格与对角抽选的证明。
\item\label{b1:item:12-holder-measures}
设 \(\mathcal P\) 为 \(E\) 上的 Borel 概率测度族，并且对某个 \(r>0\)，
\[
 \sup_{\rho\in\mathcal P}\int_E\bigl(|f(0)|+H_\alpha(f)\bigr)^r\,\mathrm d\rho(f)<\infty.
\]
说明被积函数的可测性，并证明 \(\mathcal P\) 相对弱紧。
\item\label{b1:item:12-holder-insufficient}
证明若只假设 \(\sup_{\rho\in\mathcal P}\int\|f\|_\infty^r\,\mathrm d\rho(f)<\infty\)，相对弱紧性可能失败。
\end{enumerate}
\end{exercise}
% 来源：自拟；把紧性应用于具体连续函数空间，不使用未建立的 Kolmogorov
% 连续性定理或 Arzela--Ascoli 定理；所需紧集在解答内实际构造。
\begin{solution}
若 \(f\in K_M\)，则 \(\|f\|_\infty\le2M\)。给定 \(K_M\) 中的序列 \((f_n)\)，在 \([0,1]\) 的可数稠密有理点集上逐点使用有界实数列的子列定理，再取对角子列 \((f_{n_j})\)，使其在每个有理点的值都收敛。

给定 \(\varepsilon>0\)，取足够细的有限有理网格，使每个 \(t\in[0,1]\) 都距某个网格点 \(t_i\) 小于 \(\delta\)，其中 \(2M\delta^\alpha<\varepsilon/2\)。由于网格有限，对充分大的 \(j,k\)，各网格点都满足 \(|f_{n_j}(t_i)-f_{n_k}(t_i)|<\varepsilon/2\)。于是
\[
 |f_{n_j}(t)-f_{n_k}(t)|
 \le2M|t-t_i|^\alpha+|f_{n_j}(t_i)-f_{n_k}(t_i)|<\varepsilon.
\]
所以该子列一致 Cauchy，从而一致收敛到连续函数 \(f\)。令 \(j\to\infty\)，定义 \(K_M\) 的每条不等式都保留，故 \(f\in K_M\)。这证明序列紧性；度量空间中的序列紧性等价于紧性，第一问得证。

在 \(H_\alpha\) 的上确界中只取不同的有理点 \(s,t\)，数值不变，因为任意固定的 \(s<t\) 都能用这样的点对逼近。每个差商都是 \(E\) 上的连续函数，故可数上确界 \(H_\alpha\) 为 Borel 函数。若第二问的共同积分上界为 \(C\)，则
\[
 \rho(E\setminus K_M)
 \le\rho\{\,|f(0)|+H_\alpha(f)>M\,\}\le\frac C{M^r}.
\]
由第一问，\(\mathcal P\) 为紧族。空间 \(E\) 完备，且有理网格上取有理值的折线函数组成可数稠密集：先由一致连续性用细网格插值，再把有限个顶点值作有理逼近即可。因此 \(E\) 是波兰空间，\cref{b1:thm:prokhorov}给出相对弱紧性。

最后令 \(f_n(t)=\sin(2\pi nt)\)、\(\rho_n=\delta_{f_n}\)。其一致范数矩恒为 \(1\)。但若 \(n\ne m\)，三角函数的正交积分给出
\[
 \int_0^1|f_n(t)-f_m(t)|^2\,\mathrm dt=1,
 \quad\|f_n-f_m\|_\infty\ge1.
\]
所以没有紧集能包含所有 \(f_n\)。若 \(\{\rho_n\}\) 为紧族，取误差 \(1/2\)，便应有紧集包含每个 \(f_n\)，矛盾。因此这一概率测度族不紧，由 Prokhorov 定理也不相对弱紧。
\end{solution}

\begin{exercise}\label{b1:ex:12-joint-marginals}
设 \(S,T\) 为波兰空间，\(\mu,\nu\) 分别为其上的概率测度。令 \(\Pi(\mu,\nu)\) 为 \(S\times T\) 上第一、第二边缘测度分别为 \(\mu,\nu\) 的概率测度全体。
\begin{enumerate}
\item\label{b1:item:12-couplings-compact}
证明 \(\Pi(\mu,\nu)\) 非空、凸且弱紧。
\item\label{b1:item:12-couplings-not-unique}
在 \(S=T=[0,1]\) 上构造一列联合概率测度，使两个边缘始终为均匀概率测度，但联合测度本身不弱收敛。用一个明确的连续函数辨认两个不同的子列极限。
\end{enumerate}
\end{exercise}
% 来源：自拟；与乘积测度题配对，说明独立乘积结构是额外信息。
% 联合测度族的紧性并不蕴涵指定序列只有一个极限。
\begin{solution}
乘积测度 \(\mu\otimes\nu\) 属于 \(\Pi(\mu,\nu)\)，所以该集合非空；推前的线性性保证其凸性。波兰空间上的概率测度为 Radon 测度，故对 \(\varepsilon>0\)，可取紧集 \(K\subset S\)、\(L\subset T\)，使 \(\mu(S\setminus K)<\varepsilon/2\)、\(\nu(T\setminus L)<\varepsilon/2\)。每个 \(\lambda\in\Pi(\mu,\nu)\) 都满足
\[
 \lambda\bigl((S\times T)\setminus(K\times L)\bigr)
 \le\mu(S\setminus K)+\nu(T\setminus L)<\varepsilon.
\]
所以 \(\Pi(\mu,\nu)\) 为紧族，因而相对弱紧。

还须证明它弱闭。第一边缘为 \(\mu\)，等价于对每个 \(f\in C_b(S)\) 有
\[
 \int_{S\times T}f(x)\,\mathrm d\lambda(x,y)=\int_Sf\,\mathrm d\mu;
\]
有界连续函数能够唯一确定概率测度，故这组等式确实刻画第一边缘。每条等式定义一个弱闭集，第二边缘也有同样的刻画。因此 \(\Pi(\mu,\nu)\) 为弱闭集，结合相对弱紧性得到弱紧性。

令 \(m\) 为 \([0,1]\) 上的均匀概率测度，定义连续映射
\[
 A(t)=(t,t),\quad B(t)=(t,1-t).
\]
令 \(\lambda_{2n}=A_\#m\)、\(\lambda_{2n-1}=B_\#m\)。两种测度的两个边缘都是 \(m\)，但取 \(h(x,y)=xy\)，有
\[
 \int h\,\mathrm d\lambda_{2n}=\int_0^1t^2\,\mathrm dt=\frac13,
 \quad
 \int h\,\mathrm d\lambda_{2n-1}=\int_0^1t(1-t)\,\mathrm dt=\frac16.
\]
因此全列不弱收敛；偶数和奇数子列分别恒等于两个不同的测度。相同边缘并未规定两个坐标之间的联合关系。
\end{solution}
